% GENERATED FILE. Edit canonical lessons, then run npm run generate:books.
% canonical-source-hash: fnv1a64:4b7d98c83f9f992f
% canonical-lessons: KA-C212-bandaru, KA-C212-karkhane, KA-C212-citramandira, KA-C212-maidana, KA-C212-dikku

\chapter{Port, Factory, Cinema hall, Ground, Direction}
\label{ch:ka-port-factory-cinema-hall-ground-direction}
\hlchaptermodality{212}

\begin{chapteropening}
\textbf{By the end of this chapter:} \emph{I can say a port, a factory, a cinema hall, a ground and a direction.}

Complete the last lesson of chapter 212: I can say a port, a factory, a cinema hall, a ground and a direction.
\end{chapteropening}

\section[\texorpdfstring{\kn{ಬಂದರು} (bandaru)}{ಬಂದರು (bandaru)}]{\kn{ಬಂದರು} (bandaru) — a port, a harbour}

\label{lesson:KA-C212-bandaru}



\begin{quote}
Before the new one: say the Kannada for for now, then the Kannada for a period.
\end{quote}

\subsection*{You'll want to know: \kn{ಬಂದರು}}
\textbf{\kn{ಬಂದರು}} — \emph{bandaru} — \textquotedblleft{}a port, a harbour\textquotedblright{}.

\begin{grammarlens}[title={What you've built}]
One more word for this chapter.
\end{grammarlens}

\subsection*{Your turn}
\begin{itemize}
\raggedright
  \item \emph{Say it:} \emph{bandaru}
  \item \emph{Say it:} \emph{bandaru}, once more
  \item \emph{Say it:} say \emph{bandaru}
  \item \emph{Recall:} say the Kannada for for now, then the Kannada for a period, then say \emph{bandaru} again
\end{itemize}

\begin{tcolorbox}[breakable,colback=teal!4,colframe=teal!35!black,title={Before you move on}]
What does \kn{ಬಂದರು} mean? (\textquotedblleft{}A port, a harbour\textquotedblright{}.) Say it once more.
\end{tcolorbox}
\section[\texorpdfstring{\kn{ಕಾರ್ಖಾನೆ} (kārkhāne)}{ಕಾರ್ಖಾನೆ (kārkhāne)}]{\kn{ಕಾರ್ಖಾನೆ} (kārkhāne) — a factory}

\label{lesson:KA-C212-karkhane}



\begin{quote}
Before the new one: say the Kannada for a period, then the Kannada for a port.
\end{quote}

\subsection*{You'll want to know: \kn{ಕಾರ್ಖಾನೆ}}
\textbf{\kn{ಕಾರ್ಖಾನೆ}} — \emph{kārkhāne} — \textquotedblleft{}a factory\textquotedblright{}.

\begin{grammarlens}[title={What you've built}]
One more word for this chapter.
\end{grammarlens}

\subsection*{Your turn}
\begin{itemize}
\raggedright
  \item \emph{Say it:} \emph{kārkhāne}
  \item \emph{Say it:} \emph{kārkhāne}, once more
  \item \emph{Say it:} say \emph{kārkhāne}
  \item \emph{Recall:} say the Kannada for a period, then the Kannada for a port, then say \emph{kārkhāne} again
\end{itemize}

\begin{tcolorbox}[breakable,colback=teal!4,colframe=teal!35!black,title={Before you move on}]
What does \kn{ಕಾರ್ಖಾನೆ} mean? (\textquotedblleft{}A factory\textquotedblright{}.) Say it once more.
\end{tcolorbox}
\section[\texorpdfstring{\kn{ಚಿತ್ರಮಂದಿರ} (citramandira)}{ಚಿತ್ರಮಂದಿರ (citramandira)}]{\kn{ಚಿತ್ರಮಂದಿರ} (citramandira) — a cinema hall}

\label{lesson:KA-C212-citramandira}



\begin{quote}
Before the new one: say the Kannada for a port, then the Kannada for a factory.
\end{quote}

\subsection*{You'll want to know: \kn{ಚಿತ್ರಮಂದಿರ}}
\textbf{\kn{ಚಿತ್ರಮಂದಿರ}} — \emph{citramandira} — \textquotedblleft{}a cinema hall\textquotedblright{}.

\begin{grammarlens}[title={What you've built}]
One more word for this chapter.
\end{grammarlens}

\subsection*{Your turn}
\begin{itemize}
\raggedright
  \item \emph{Say it:} \emph{citramandira}
  \item \emph{Say it:} \emph{citramandira}, once more
  \item \emph{Say it:} say \emph{citramandira}
  \item \emph{Recall:} say the Kannada for a port, then the Kannada for a factory, then say \emph{citramandira} again
\end{itemize}

\begin{tcolorbox}[breakable,colback=teal!4,colframe=teal!35!black,title={Before you move on}]
What does \kn{ಚಿತ್ರಮಂದಿರ} mean? (\textquotedblleft{}A cinema hall\textquotedblright{}.) Say it once more.
\end{tcolorbox}
\section[\texorpdfstring{\kn{ಮೈದಾನ} (maidāna)}{ಮೈದಾನ (maidāna)}]{\kn{ಮೈದಾನ} (maidāna) — a ground, a playground}

\label{lesson:KA-C212-maidana}



\begin{quote}
Before the new one: say the Kannada for a factory, then the Kannada for a cinema hall.
\end{quote}

\subsection*{You'll want to know: \kn{ಮೈದಾನ}}
\textbf{\kn{ಮೈದಾನ}} — \emph{maidāna} — \textquotedblleft{}a ground, a playground\textquotedblright{}.

\begin{grammarlens}[title={What you've built}]
One more word for this chapter.
\end{grammarlens}

\subsection*{Your turn}
\begin{itemize}
\raggedright
  \item \emph{Say it:} \emph{maidāna}
  \item \emph{Say it:} \emph{maidāna}, once more
  \item \emph{Say it:} say \emph{maidāna}
  \item \emph{Recall:} say the Kannada for a factory, then the Kannada for a cinema hall, then say \emph{maidāna} again
\end{itemize}

\begin{tcolorbox}[breakable,colback=teal!4,colframe=teal!35!black,title={Before you move on}]
What does \kn{ಮೈದಾನ} mean? (\textquotedblleft{}A ground, a playground\textquotedblright{}.) Say it once more.
\end{tcolorbox}
\section[\texorpdfstring{\kn{ದಿಕ್ಕು} (dikku)}{ದಿಕ್ಕು (dikku)}]{\kn{ದಿಕ್ಕು} (dikku) — a direction}

\label{lesson:KA-C212-dikku}



\begin{quote}
Before the new one: say the Kannada for a cinema hall, then the Kannada for a ground.
\end{quote}

\subsection*{You'll want to know: \kn{ದಿಕ್ಕು}}
\textbf{\kn{ದಿಕ್ಕು}} — \emph{dikku} — \textquotedblleft{}a direction\textquotedblright{}.

\begin{grammarlens}[title={What you've built}]
One more word for this chapter.
\end{grammarlens}

\subsection*{Your turn}
\begin{itemize}
\raggedright
  \item \emph{Say it:} \emph{dikku}
  \item \emph{Say it:} \emph{dikku}, once more
  \item \emph{Say it:} say \emph{dikku}
  \item \emph{Recall:} say the Kannada for a cinema hall, then the Kannada for a ground, then say \emph{dikku} again
\end{itemize}

\begin{tcolorbox}[breakable,colback=teal!4,colframe=teal!35!black,title={Before you move on}]
What does \kn{ದಿಕ್ಕು} mean? (\textquotedblleft{}A direction\textquotedblright{}.) Say it once more.
\end{tcolorbox}
